{\bfseries

Deutsch-Jozsa con la función XOR. Considere la función $f:\{0,1\}^2 \rightarrow \{0,1\}$ definida por $f(00)=0$, $f(01)=1$, $f(10)=1$, $f(11)=0$ (es decir, la función XOR de los dos bits).
Aplique el algoritmo de Deutsch-Jozsa paso a paso para $n=2$, mostrando el estado después de cada etapa:
\begin{itemize}
    \item Inicialización: $\vert \psi_0 \rangle = \vert 00 \rangle \otimes \vert 1 \rangle$.
    \item Primera Hadamard sobre el primer registro.
    \item Aplicación del oráculo $U_f$ (que actúa como $\vert x \rangle \vert y \rangle \mapsto \vert x \rangle \vert y \oplus f(x) \rangle$).
    \item Segunda Hadamard sobre el primer registro.
\end{itemize}
Finalmente, calcule la probabilidad de medir $\vert 00 \rangle$ en el primer registro y verifique
que es nula, lo que confirma que la función es balanceada.}

\vspace{.5cm}

\begin{quote}
    \begin{enumerate}
        \item \textbf{Inicialización }
            $$
            \vert \psi_0 \rangle = \vert 00 \rangle \otimes \vert 1 \rangle
            $$

        \item \textbf{Primera Hadamard sobre el primer registro.* } (aquí dice solo sobre
            el primer registro pero cuando busque para hacerlo todos aplican H a todos los
            registros, seguí esta práctica)

            \begin{align*}
                \vert \psi_1 \rangle = H^{\otimes 3} \vert \psi_0 \rangle  &=  H^{\otimes 2}\vert 00
                \rangle  H\otimes \vert 1 \rangle  \\ 
                                                                            &=  \frac{1}{2\sqrt{2}}
                                                                            (\ket{0}+\ket{1})
                                                                            \otimes (\ket{0}+\ket{1})
                                                                            \otimes (\ket{0}-\ket{1})\\
            \end{align*}

        \item \textbf{Aplicación del oráculo $U_f$}

            Primero voy a reescribir para que se vea mas bonito:
            \begin{align*}
                \frac{1}{2\sqrt{2}} (\ket{0}+\ket{1}) \otimes (\ket{0}+\ket{1}) \otimes
                (\ket{0}-\ket{1}) &= \frac{1}{2} (\ket{00} + \ket{01} + \ket{10} + \ket{11})
                \otimes \ket{-} \\
                                  &= \frac{1}{2} (\ket{00}\ket{-} + \ket{01}\ket{-} + \ket{10}\ket{-} + \ket{11}\ket{-}) 
            \end{align*}

            Como nuestro segundo registro esta en $\ket{-}$, sabemos que habra phase kickback y
            podemos reducir la regla:
            $$
            U_f \ket{x} \ket{-} = (-1)^{f(x)}  \ket{x} \ket{-}
            $$

            Con esto, y el los valores de la función resolvemos:

            \begin{align*}
                \ket{\psi_2} = \frac{1}{2} U_f(\ket{00}\ket{-} + \ket{01}\ket{-} + \ket{10}\ket{-} +
                \ket{11}\ket{-}) = \\ 
                \frac{1}{2} (U_f\ket{00}\ket{-} + U_f\ket{01}\ket{-} + U_f\ket{10}\ket{-} +
                U_f\ket{11}\ket{-}) = \\
                \frac{1}{2} (\ket{00}\ket{-} -\ket{01}\ket{-} -\ket{10}\ket{-}
                + \ket{11}\ket{-})
            \end{align*}

        \item \textbf{Segunda Hadamard sobre el primer registro.}

            Vamos a reescribir para que quede mas bonito:
            $$
            \ket{\psi_2} = \frac{1}{2} (\ket{0} (\ket{0} -\ket{1})
            -\ket{1} (\ket{0} - \ket{1})) \otimes \ket{-} = \frac{1}{2} ((\ket{0}
            -\ket{1}) \otimes (\ket{0} - \ket{1})) \otimes \ket{-} =
            \ket{-}\otimes \ket{-}\otimes \ket{-}
            $$

            Ahora, trivialmente:
            \begin{align*}
                \ket{\psi_3} = H^{\otimes 2} \ket{\psi_2} = H^{\otimes
                2}(\ket{-}\otimes \ket{-}\otimes \ket{-}) &= H \ket{-}\otimes
                H \ket{-} \otimes \ket{-} \\
                                                          &= \ket{1} \otimes \ket{1} \otimes \ket{-} \\
            \end{align*}
    \end{enumerate}

    \textbf{Calcular la probabilidad de medir $\ket{00}$ en el primer registro y confirmar que la
    función es balanceada.}

    $$
    P(00) = |\bra{00}\ket{11}| ^2 = 0
    $$

    Sabemos que es igual a 0 porque son ortonormales. Debido a que la probabilidad es exactamente 0
    se verifica que la función es balanceada.
\end{quote}

\vspace{.5cm}
